\cvheader{Arsenij Schützer}{Фуллстек-разработчик}{%
  \begin{tabular}{@{}l@{\hspace{0.6em}}l@{}}
    {\color{dim}Адрес:}   & Königsberger Str.~4B, 75181 Pforzheim, Германия\\
    {\color{dim}Родился:} & 08.10.1992, Омск (Россия)\\
    {\color{dim}Телефон:} & +49 162 629 1630\\
    {\color{dim}E-Mail:}  & \href{mailto:senyaarseniymail@gmail.com}{senyaarseniymail@gmail.com}\\
    {\color{dim}Web:}     & \href{https://senyaak.work}{senyaak.work} \textcolor{dim}{·} \href{https://github.com/senyaak}{github.com/senyaak} \textcolor{dim}{·} \href{https://linkedin.com/in/senyaak}{linkedin.com/in/senyaak}
  \end{tabular}
}

\section*{О себе}

Фуллстек-разработчик с 8 годами коммерческого опыта. В профессию пришёл через самостоятельно начатое дистанционное обучение информатике (ils.de, 2014/15) и трёхлетнюю профессиональную подготовку по специальности Fachinformatiker Anwendungsentwicklung (разработчик приложений, 2015–2018). Основной стек: Angular, Vue, TypeScript, Node.js, Java, PostgreSQL. Опыт самостоятельной реализации приложений целиком (фронтенд, бэкенд, база данных, архитектура), а также работы в кросс-функциональных командах. Родной язык русский.

\section*{Технические навыки}

\cvkvblock{
\cvkv{Основной стек}{Angular (все версии начиная с AngularJS), Vue~3, TypeScript/JavaScript, Node.js, Java, PostgreSQL, MongoDB}
\cvkv{Уверенно}{React, NestJS, Python (FastAPI, Kedro), Linux, Git, Docker, WebSocket, REST API, JSON Schema}
\cvkv{Также работал с}{C, C++, C\#, PHP, Rust, Unity, Jenkins, Kubernetes, Socket.IO, LLM Agents}
}

\section*{Опыт работы}

\cvjob{e.telligent GmbH, Regensburg, Германия}{09/2023 – сейчас}{Фуллстек-разработчик}{e.telligent разрабатывает программные решения для электромобильности. Работа над продуктом e.CoSys, ASPICE-совместимой платформой управления жизненным циклом приложений (ALM) для автомобильной разработки со сквозной трассируемостью по V-модели.}
\begin{itemize}
  \item Фуллстек-разработка по всему стеку: \textbf{Vue~3 + TypeScript} на фронтенде, \textbf{Java} с \textbf{PostgreSQL} на бэкенде
  \item Проектирование и реализация функций под конкретных заказчиков на обобщённой лицензируемой архитектуре платформы
  \item Самостоятельная разработка сложных фронтенд-компонентов
  \item Бэкенд-разработка с упором на сложные структуры данных
  \item Код-ревью и refinement-сессии в команде
  \item Сервис-обёртка для ИИ на Python
\end{itemize}

\cvjob{Mittwald CM Service GmbH \& Co.~KG, Espelkamp, Германия}{07/2022 – 12/2022}{Фуллстек-разработчик}{Mittwald: крупный немецкий хостинг-провайдер. Работа над клиентским порталом в составе большой команды разработки.}
\begin{itemize}
  \item Фронтенд-разработка клиентского портала на \textbf{React} и \textbf{TypeScript}
  \item Работа в agile-команде с устоявшимися стандартами кода и процессами CI/CD
\end{itemize}

\cvjob{saltation GmbH \& Co.~KG, Bielefeld, Германия}{08/2018 – 01/2022}{Фуллстек-разработчик}{saltation разрабатывает индивидуальное ПО для разных отраслей (20–40 сотрудников). Самостоятельная реализация нескольких проектов от фронтенда до базы данных, в отдельных случаях в тандеме (фронтенд в одиночку плюс бэкендер).}
\begin{itemize}
  \item CMS для ведения данных об объектах (порты, велосипеды): обобщённые модели на основе JSON Schema с автоматической генерацией форм через конфигурацию JSON Form. Стек: \textbf{AngularJS, Node.js, PostgreSQL}; единоличная ответственность на протяжении нескольких лет
  \item Архитектурное расширение CMS для интернационализации (параллельный многоязычный контент): согласованные изменения в базе данных, бэкенде и фронтенде
  \item Система моделирования бизнес-процессов (обучение, повышение квалификации): \textbf{AngularJS} с собственным движком, несколько вариантов фронтенда
  \item Инструмент администрирования проката оборудования: \textbf{Angular} (начиная с версии~2)
  \item Приложение для онбординга на планшетах на \textbf{Unity (C\#)}
  \item Эпизодическая бэкенд-работа на \textbf{Rust} (доработки на сервере), а также небольшие проекты на Java и C\#
\end{itemize}

\cvjob{saltation GmbH \& Co.~KG, Bielefeld, Германия}{08/2015 – 07/2018}{Профессиональная подготовка: Fachinformatiker Anwendungsentwicklung (IHK)}{}
\begin{itemize}
  \item Выпускной проект: проектирование и полная реализация расширения интернационализации для существующей CMS (база данных, бэкенд, фронтенд; в одиночку)
  \item Практика в нескольких веб-проектах на AngularJS, Node.js, PostgreSQL
\end{itemize}

\section*{Личные проекты}

\begin{itemize}
  \item \textbf{VTT (Virtual Tabletop):} собственная система виртуального стола для настольных ролевых игр (аналог Foundry/Roll20). Бэкенд на Node.js + WebSocket (сервер как единый источник правды), рендеринг на Canvas, механики зон и поверхностей, правила в JSON
  \item \textbf{Comersant, многопользовательская настольная игра:} многолетний личный проект (онлайн-игра в духе «Монополии») на Angular, Node.js, WebSocket. Архитектурные решения задокументированы в ADR
  \item \textbf{hochstrasser-immobilien.de:} сайт на безвозмездной основе, включая SEO-оптимизированный немецкий контент, для агентства недвижимости в Ульме/Ной-Ульме
  \item \textbf{senyaak.work:} личный сайт на трёх языках (DE/EN/RU)
\end{itemize}

\section*{Языки}

\cvkvblock{
\cvkv{Русский}{родной}
\cvkv{Немецкий}{свободно (C1), в Германии с 2009 года; профессиональная подготовка и вся карьера на немецком}
\cvkv{Английский}{очень хорошо устно и письменно (B2/C1), регулярно читаю техническую и общую литературу}
}

\section*{Профессиональное обучение}

\cventry{08/2015 – 07/2018}{Fachinformatiker Anwendungsentwicklung (разработчик приложений, IHK)}{saltation GmbH \& Co.~KG, Bielefeld, Германия}
\cventry{02/2014 – 04/2015}{Дистанционное обучение по информатике}{ils.de}

\section*{Образование}

\cventry{09/2013 – 06/2015}{Гимназия CJD}{Versmold, Германия}
\cventry{09/2011 – 06/2013}{Аттестат Realschule, Akademie Klausenhof}{Rhede, Германия}
\cventry{09/2009 – 06/2011}{Аттестат Hauptschule, WRSD am Deutenberg}{VS-Schwenningen, Германия}
\cventry{09/2001 – 06/2008}{Средняя общеобразовательная школа №14}{Омск, Россия}
